\section{A probe-rs session, and the three traps in front of it}
\label{sec:a-probe-rs-session}

\texttt{probe-rs} does four things in one command on this board: it erases and flashes the image, resets the core, attaches
a debugger, and prints whatever the firmware logs through RTT. On a board with no UART broken out and no pins to spare,
that last part is the difference between a program you can watch and a program you can only guess at.

\begin{minted}{bash}
probe-rs run --chip EFR32MG24B220F1536IM48 \
  target/thumbv8m.main-none-eabihf/release/examples/hal_rtt_hello
\end{minted}

What came back, on the board, with \texttt{defmt} and a heartbeat loop:

\begin{minted}{text}
Finished in 0.47s
1789052999 [INFO ] rtt: heartbeat 0 -- 78 MHz (hal_rtt_hello .../hal_rtt_hello.rs:44)
1797056423 [INFO ] rtt: heartbeat 1 -- 78 MHz (hal_rtt_hello .../hal_rtt_hello.rs:44)
\end{minted}

The number is the DWT cycle counter, the level is the record's, and the location is where the record was written --- all of
it decoded on the host, because with deferred formatting the strings never travel.

\subsection{The three traps, in the order they bit}

\textbf{1. \texttt{-Tdefmt.x}, or the probe finds a version and no section.} The first run failed with
\texttt{defmt version found, but no `.defmt` section --- check your linker configuration}: the section that carries the
format strings has to be placed by the \texttt{defmt.x} linker script. Putting it in \texttt{.cargo/config.toml} fixes
one crate and breaks every other one, because their examples have no \texttt{defmt} build script to supply the file; it
belongs in the logging crate's \texttt{build.rs}, as
\texttt{cargo:rustc-link-arg-examples=-Tdefmt.x}, where it reaches that crate's examples and nothing else.

\textbf{2. The panic handler wants a timestamp.} \texttt{panic-probe} prints panics through \texttt{defmt}, and
\texttt{defmt} refuses to link without \texttt{\_defmt\_timestamp}. The DWT cycle counter is the cheapest real clock
here, and the way to supply it is a macro invocation, not an attribute:

\begin{minted}{Rust}
defmt::timestamp!("{=u64}", unsafe {
    (*cortex_m::peripheral::DWT::PTR).cyccnt.read() as u64
});
\end{minted}

The attribute form does not exist in this version, and the counter has to be enabled in \texttt{main} before the first
record.

\textbf{3. \texttt{DEFMT\_LOG} defaults to \texttt{error}.} With the log level unset, \texttt{info!} and \texttt{warn!}
are compiled out of the binary entirely, so the board runs correctly, sends nothing, and every layer looks healthy. This
one is worth knowing because the target side is verifiably fine while it happens --- reading the RTT control block out of
RAM showed the id string in place, one up channel, \texttt{size = 1024}, both pointers valid, and \texttt{write = 0}:

\begin{minted}{text}
20000008: 47474553 52205245 00005454 00000000   # "SEGGER RTT"
20000018: 00000001 00000000 20000000 20000044   # max_up, max_down, name, buffer
20000028: 00000400 00000000 00000000 00000001   # size, wr=0, rd=0, flags
\end{minted}

A control block that is perfectly formed and a write index that never moves: the target is initialised and has nothing to
say, which is a different fault from a broken transport. \texttt{DEFMT\_LOG = "info"} in \texttt{[env]} is the fix.

\subsection{A fourth thing that looks like a fault}

The up channel is non-blocking by design: a record written before the probe has attached is dropped, because there is
nobody to hand it to. A program that logs once at boot and then sleeps therefore looks silent no matter how correct it
is. Log a heartbeat instead --- and if a fault kills the program before the probe drains the buffer, that is the case
where RTT tells you nothing and a debugger does.

\subsection{When the target will not halt}

\texttt{openocd} reported \texttt{target was in unknown state when halt was requested} on the same board in the same
hour, so the memory read had to come from the probe instead:

\begin{minted}{bash}
probe-rs read b32 0x20000008 12 --chip EFR32MG24B220F1536IM48
\end{minted}

Reading memory needs no halt, which makes it the tool of choice for inspecting a running target's own structures --- the
RTT block above, a mailbox, a ring buffer --- without disturbing what it is doing.
